\section{开发组织：让反例进入设计}
\label{sec:process}

本项目将开发组织视为系统设计的一部分：并行工作的边界必须与状态所有权一致，而验证必须能够推翻当前实现。最初的共享契约先固定帧格式、认证身份、提交后发布、逐会话先进先出等不变量，再划分协议与套接字、服务端存储、客户端及集成责任。这样，工作者交换的是明确的接口与约束，而不是对彼此代码的隐含猜测。终端重构沿用同一原则：依赖与原生输入机制、纯视图构造、会话传输与应用状态、真实终端测量分别归属不同任务；涉及接口的变更先回到共享契约，再进入集成。分工不是尽可能多地同时修改文件，而是减少需要共同理解和共同修改的状态。

\paragraph{文档是可修订的共享黑板。}
架构契约、存储语义、终端决策与验证记录构成持久的共享黑板（shared blackboard）。例如，终端方案从自建输入与绘制引擎转向固定版本的 FTXUI 后，旧草案保留为决策轨迹，却明确退出当前实现要求；否则，历史方案容易成为并行任务之间相互矛盾的隐性规范。评审记录同样区分初始发现、后续修正与执行证据：一条已修复的历史缺陷仍有解释价值，但不能被继续当作当前结论。黑板由此承担两项职责：使后来者恢复设计理由，并使团队能够撤销已经失效的认识。

\paragraph{独立评审寻找跨边界反例。}
最有价值的缺陷并非局部语法错误，而是两个分别合理的模块组合后破坏了整体语义。R1 发现，发起连接关闭后丢弃整个工作完成项，会同时丢弃已经提交、应向其他连接广播的消息；修复将请求者生命周期与领域事件发布解耦。R2 则构造了另一种调度：较小编号已经提交，但其工作线程尚未入队，较大编号先发布；接收端随后提交较大已读游标，便可能将尚未显示的消息标记为已读。这个反例说明，逐会话顺序与数据库内的房间锁不足以保证跨发送者的显示顺序，房间级顺序边界必须延伸到完成项入队。评审的产物因此不只是问题清单，也是对所有权边界的重新刻画。

客户端评审继续沿着数据到用户判断的路径追踪语义。T1 揭示房间最大编号、实时消息编号和连续取回的历史游标不能共用一个变量；元数据声称存在某条消息，不等于客户端已经取得此前的历史。T4 发现，已转发但未收到确认的发送被表示为确定失败，会诱导用户用新随机标识重发；交付不确定性必须跨守护进程与界面保持一致。T5 则暴露了有界缓存与历史浏览的冲突：若始终淘汰最旧消息，刚取回的旧历史会立即消失。改为历史窗口后，还需处理返回最新时旧分页仍在途的取消竞态；后续响应必须服从当前浏览意图，而不能凭旧请求上下文重新启用历史淘汰策略。

\paragraph{用户旅程决定证据是否相关。}
用户报告 Windows 客户端逐字迟滞且对话难以阅读时，既有管道测试的成功并不能反驳这一观察：它们绕过了原生键盘、终端布局与输出刷新。验证因而转向真实伪终端（pseudoterminal, PTY）与 Windows 伪控制台（ConPTY），通过独立协议端点确认实际发送，并检查当前屏幕、编辑中的草稿、重连和终端恢复，而非仅搜索累计输出。性质测试（property-based testing）用于探索字节流分片、UTF-8 和游标状态空间\cite{quickcheck}；确定性旅程则固定跨组件反例。历史回归先让旧二进制停在第 200 条而非第 300 条，再验证修复；取消竞态则扣住旧页结束响应，执行返回最新，随后释放响应，检查旧请求不会复活。这样的失败具有设计指向性，而非一次笼统超时。

\paragraph{测试设施本身也必须接受反例。}
macOS 执行暴露了另一个边界：测试线程等待协议端点时不读取 PTY，可能使同步屏幕输出阻塞。修正是持续排空输出，而非修改应用输入解析。随后，控制会话退出后的终端属性检查改为由仍存活的会话启动器在挂断前记录；对于内核自动设置的输入重处理状态，则在前后快照中一致地执行不消耗输入的就绪查询，仍要求完整属性相等。负向探针保留错误的回显设置，确认恢复检查依然能够拒绝缺陷。由此得到的组织经验不是未经对照的生产率提升，而是一条可检查的闭环：共享契约约束并行实现，独立反例修订模型，真实旅程检验用户可见结果，文档保存下一轮工作所需的理由。
